\documentclass[11pt,a4paper]{article}
\usepackage[T1]{fontenc}
\usepackage[margin=25mm]{geometry}
\usepackage{amsmath,amssymb,booktabs}
\usepackage[hidelinks]{hyperref}
\usepackage{microtype}

\newcommand{\dd}{\mathrm{d}}
\newcommand{\ii}{\mathrm{i}}
\newcommand{\code}[1]{\texttt{\detokenize{#1}}}
\DeclareMathOperator{\sgn}{sgn}
\DeclareMathOperator{\atanTwo}{atan2}
\DeclareMathOperator{\diag}{diag}

\title{Response matrices for the FCC-ee correction workflow}
\author{FCC-ee Xsuite tuning}
\date{25 October 2026}

\begin{document}
\maketitle

\begin{abstract}
This note describes the response matrices implemented in this repository:
closed orbit, normal optics and horizontal dispersion, coupling and vertical
dispersion, tunes, and first- and second-order chromaticities. Equations use
the same signs, coordinates, and knob units as the code. Analytical responses
are local, uncoupled-optics approximations; numerical family responses are
obtained from Twiss probes. The appendices explain the four-point
Gauss--Legendre integration and the propagation of optics inside a quadrupole.
\end{abstract}

\tableofcontents

\section{Reference machine and matrix convention}

\code{prepare_reference} in \code{correction_workflow.py} prepares a common
reference with errors off, correction trims zero, and nominal nonlinear
strengths (\code{on_nonlinear=1}). It computes periodic 4D Twiss at
\code{delta0=0} with chromatic quantities enabled. Correction is performed
without radiation or tapering. All response builders use this clean reference;
the resulting matrices remain fixed during correction, while residuals are
evaluated on the current machine.

Let $\boldsymbol{u}$ be the physical correction knobs, $F_a$ an observable,
and $T_a$ its target. The stored matrix is
\begin{equation}
 r_a(\boldsymbol{u})=F_a(\boldsymbol{u})-T_a,\qquad
 R_{aj}=\left.\frac{\partial r_a}{\partial u_j}\right|_{\rm ref},\qquad
 \boldsymbol{r}(\boldsymbol{u}+\Delta\boldsymbol{u})
 \simeq\boldsymbol{r}(\boldsymbol{u})+R\Delta\boldsymbol{u}.
 \label{eq:definition}
\end{equation}
Targets are independent of the correction knobs. Orbit targets are live BPM
offset expressions: they change with the error ramp, but not with an orbit
trim at a fixed ramp setting. Thus the stored responses are derivatives of
the observables themselves.

Rows and columns follow the optimizer's target and vary lists. The matrix
contains neither target weights nor knob scales. If $W$ contains the target
weights and $S$ the \code{Vary.weight} scales, the solver uses
\begin{equation}
 A=WRS,\qquad
 \Delta\boldsymbol{u}=-S A^+ W\boldsymbol{r}.
 \label{eq:solve}
\end{equation}
Here $A^+$ is the SVD pseudoinverse with the configured singular-value cutoff;
only active rows and columns enter the solve. Step scaling and limits may
shorten the resulting change. The workflow evaluates the changed machine
and checks its correction goals separately.

\begin{center}
\begin{tabular}{lll}
\toprule
Module & Observables & Response construction \\
\midrule
Orbit ($x$, $y$) & $x$, $y$ at BPMs & Closed-orbit Green function \\
Optics & $\beta_{x,y},\alpha_{x,y},D_x,Q_{x,y}$ & Thick-quadrupole integrals \\
Coupling & $\Re f,\Im f,D_y$ & Skew-coil midpoint model \\
Tune & $Q_x,Q_y$ & Forward Twiss probes \\
Chromaticity & $Q'_x,Q'_y$ & Central Twiss probes \\
Second order & $Q''_x,Q''_y$ & Central Twiss probes \\
\bottomrule
\end{tabular}
\end{center}

Throughout, $p\in\{x,y\}$ denotes a transverse plane,
$\sigma_x=+1$, $\sigma_y=-1$, and $Q_p$ is its tune. The phase $\mu_p(s)$
is measured in \emph{turns}, as in Twiss, so its radian value is $2\pi\mu_p$.
All optics functions in the following kernels are reference quantities.
The longitudinal coordinate is $s$; $\delta$ is the relative momentum deviation.

\section{Closed-orbit response}

For observation $s_i$ and a positive kick $\theta_{pj}$ at $s_j$, the periodic
uncoupled Green function is
\begin{equation}
 G_p(s_i,s_j)=
 \frac{\sqrt{\beta_p(s_i)\beta_p(s_j)}}{2\sin(\pi Q_p)}
 \cos\!\left[2\pi|\mu_p(s_i)-\mu_p(s_j)|-\pi Q_p\right].
 \label{eq:green}
\end{equation}
\code{orbit_response} returns $R_{ij}=G_p(s_i,s_j)$, with rows at BPMs and
columns at correctors. Knobs represent kicks in radians. Their installation
uses $\Delta\code{knl[0]}=-\theta_x$ and
$\Delta\code{ksl[0]}=+\theta_y$, so a positive knob produces a positive
canonical transverse momentum kick in either plane.

The matching residual is $x_i-d_{x,i}$ or $y_i-d_{y,i}$, where $d_{p,i}$ is
the live BPM offset. The stored matrix is in the laboratory frame and has no
cross-plane blocks. The orbit acceptance diagnostic additionally rotates the
offset-subtracted coordinates by the BPM roll; that diagnostic rotation is
not included in the analytical matrix.

\section{Normal-quadrupole response}
\label{sec:normal}

\subsection{Distributed strength and kernels}

For quadrupole $j$ of length $L_j$, let $u_j=\Delta k_{1j}$ be a uniform
normal-gradient trim in $\mathrm{m}^{-2}$. A small slice $\dd\ell$ produces
the focusing perturbation $\sigma_p u_j\dd\ell$. Write its location as
$t_j(\ell)=s_j+\ell$ and define
\begin{equation}
 d_p(i,t)=\mu_p(s_i)-\mu_p(t),\qquad
 \Phi_p(i,t)=4\pi|d_p(i,t)|-2\pi Q_p.
\end{equation}
The implementation in \code{normal_quadrupole_response} integrates the
following kernels along each magnet:
\begin{align}
 \frac{\partial\beta_p(s_i)}{\partial k_{1j}}
 &= -\frac{\sigma_p\beta_p(s_i)}{2\sin(2\pi Q_p)}
 \int_0^{L_j}\beta_p(t_j(\ell))
       \cos\Phi_p(i,t_j(\ell))\,\dd\ell,
 \label{eq:beta}\\
 \frac{\partial\alpha_p(s_i)}{\partial k_{1j}}
 &= -\frac{\sigma_p}{2\sin(2\pi Q_p)}
 \int_0^{L_j}\beta_p(t_j(\ell))
 \Big[\sgn d_p(i,t_j(\ell))\sin\Phi_p(i,t_j(\ell))
 \nonumber\\[-2pt]
 &\hspace{46mm}+\alpha_p(s_i)\cos\Phi_p(i,t_j(\ell))\Big],\dd\ell,
 \label{eq:alpha}\\
 \frac{\partial Q_p}{\partial k_{1j}}
 &=\frac{\sigma_p}{4\pi}\int_0^{L_j}\beta_p(t_j(\ell))\,\dd\ell.
 \label{eq:tuneanalytic}
\end{align}
In particular, the beta and alpha rows are derivatives of the \emph{ordinary}
Twiss functions, not of logarithms or normalized-alpha combinations. The
cosine term proportional to $\alpha_p(s_i)$ in Eq.~\eqref{eq:alpha} is needed
for the ordinary-alpha response. Relative-beta weighting is applied later:
a target weight proportional to $1/\beta_{p,\rm ref}(s_i)$ makes its weighted
residual proportional to $\beta_p/\beta_{p,\rm ref}-1$.

A normal trim gives a horizontal momentum kick
$\Delta p_x=-u_j x\dd\ell$. On the clean reference orbit, the corresponding
dispersion source is $-u_jD_x(t)\dd\ell$, giving
\begin{equation}
 \frac{\partial D_x(s_i)}{\partial k_{1j}}
 =-\int_0^{L_j}G_x(s_i,t_j(\ell))D_x(t_j(\ell))\,\dd\ell.
 \label{eq:dx}
\end{equation}
These are linearized, uncoupled, on-momentum response formulas. They do not
include a coupled-optics or nonlinear feed-down response of the dirty machine.
The integrals are evaluated using the same four internal sample points for
beta, alpha, tune, and dispersion (Appendices~\ref{app:gauss} and
\ref{app:transport}).

\subsection{From magnet derivatives to actual knobs}

The low-level matrix $B_{aj}=\partial F_a/\partial k_{1j}$ has columns for
individual thick quadrupoles. Its rows are grouped by quantity:
$\beta_x$, $\beta_y$, $\alpha_x$, $\alpha_y$, $D_x$, $Q_x$, $Q_y$.
Beta, alpha, and dispersion can have different observation lists.

\code{build_normal_quad_optics_response_matrix} in \code{corr_optics.py}
finds the quadrupoles controlled by each independent knob through xdeps.
Let $K_{1j}$ be the magnet's total integrated normal strength, obtained from
\code{get_total_knl_ksl()[0][1]}. For knob $u_b$ and its configured probe step
$h_b$, the strength slope is measured without running Twiss:
\begin{equation}
 C_{jb}\simeq\frac{K_{1j}(u_b+h_b)-K_{1j}(u_b-h_b)}{2h_b},\qquad
 R_{ab}=\sum_j\frac{B_{aj}}{L_j}C_{jb}.
 \label{eq:chain}
\end{equation}
Each knob is restored after the probes, also on failure. The factor $1/L_j$
converts a distributed-gradient response into an integrated-strength response.
This handles knobs shared by several quadrupoles, derived family expressions,
and contributions to the total strength beyond the main \code{k1} field.
Integrated changes are modeled as distributed uniformly over each magnet.
The supported magnets are unsliced thick quadrupoles with fixed positive
lengths. The matrix is reordered into the actual target list before use.

The standard optics optimizer has five rows per observation,
$(\beta_x,\beta_y,\alpha_x,\alpha_y,D_x)$, followed by $Q_x,Q_y$.
The low-level helper also returns
\begin{equation}
 s_j^{\rm quad}=
 \left[\max\!\left(\int_0^{L_j}\beta_x(t_j(\ell))\dd\ell,
                   \int_0^{L_j}\beta_y(t_j(\ell))\dd\ell\right)\right]^{-1}.
\end{equation}
These auxiliary scales are currently discarded by \code{corr_optics.py};
its standard optics varies use unit weights.

\section{Skew-quadrupole response: coupling and vertical dispersion}

\subsection{Physical groups and midpoint optics}

Each coupling knob $u_g$ is an integrated skew-quadrupole strength in
$\mathrm{m}^{-1}$, installed on one physical sextupole. If the magnet has
several pieces with lengths $L_m$, the added skew strength is
\begin{equation}
 \Delta\code{ksl[1]}_m=u_g\frac{L_m}{L_g},\qquad L_g=\sum_{m\in g}L_m.
\end{equation}
The analytical model replaces the distributed coil by one kick at the physical
magnet midpoint $s_g=s_{g,0}+L_g/2$. It uses the entrance of the first piece and
the total length, assuming the pieces describe one contiguous magnet.
For $\ell_g=L_g/2$, midpoint optics are approximated by drift transport:
\begin{align}
 \beta_{p,g}&=\beta_{p,0}-2\alpha_{p,0}\ell_g+
                 \frac{1+\alpha_{p,0}^2}{\beta_{p,0}}\ell_g^2,\\
 \mu_{p,g}&=\mu_{p,0}+
   \frac{1}{2\pi}\atanTwo(\ell_g,\beta_{p,0}-\alpha_{p,0}\ell_g),\\
 D_{x,g}&=D_{x,0}+\ell_g D_{p_x,0}.
\end{align}
Here $D_{p_x}$ denotes \code{dpx}. This is a clean-orbit, straight-sextupole
approximation; unlike the normal-quadrupole response, it uses no Gauss points.

\subsection{Resonance-driving terms}

For an observation $s_i$, define the forward phase from source to observation
with the ring wrap chosen by longitudinal ordering:
\begin{equation}
 \Delta\mu_p(i,g)=\mu_p(s_i)-\mu_{p,g}
                  +Q_p\,\mathbf{1}_{s_g>s_i}.
\end{equation}
The complex coupling derivatives implemented in \code{skew_quadrupole_response}
are
\begin{align}
 \frac{\partial f_{1001}(s_i)}{\partial u_g}
 &=\frac{\sqrt{\beta_{x,g}\beta_{y,g}}
    e^{2\pi\ii[\Delta\mu_x(i,g)-\Delta\mu_y(i,g)]}}
          {4[1-e^{2\pi\ii(Q_x-Q_y)}]},\\
 \frac{\partial f_{1010}(s_i)}{\partial u_g}
 &=\frac{\sqrt{\beta_{x,g}\beta_{y,g}}
    e^{2\pi\ii[\Delta\mu_x(i,g)+\Delta\mu_y(i,g)]}}
          {4[1-e^{2\pi\ii(Q_x+Q_y)}]}.
 \label{eq:rdt}
\end{align}
Their real and imaginary parts form four real response blocks. These signs
and phases follow the code's positive \code{ksl[1]} and Twiss RDT conventions.

A positive integrated skew trim produces $\Delta p_y=u_g x$, so its vertical
dispersion response is
\begin{equation}
 \frac{\partial D_y(s_i)}{\partial u_g}
 =\frac{\sqrt{\beta_y(s_i)\beta_{y,g}}}
             {2\sin(\pi Q_y)}
   \cos[2\pi\Delta\mu_y(i,g)-\pi Q_y]\,D_{x,g}.
\end{equation}
The helper initially groups rows by quantity; \code{corr_coupling.py} reorders
them into five rows per observation:
$(\Re f_{1001},\Im f_{1001},\Re f_{1010},\Im f_{1010},D_y)$.
Targets are the corresponding reference values. Column scales are
$s_g^{\rm skew}=1/\sqrt{\beta_{x,g}\beta_{y,g}}$ and are supplied through
\code{Vary.weight}, rather than multiplied into the stored matrix.

During construction the code checks one distributed probe,
$u_g=10^{-5}s_g^{\rm skew}\sin(g+0.3)$ with zero-based group index $g$.
It compares the Twiss change with $R\boldsymbol{u}$, both for the full vector
and for its $D_y$ block, and rejects relative norm discrepancies above $5\%$.
The clean zero trims are restored afterwards. This directional check tests
the combined response, not every column independently.

\section{Numerical family responses}

These matrices are computed by perturbing family trims and evaluating Twiss
on the clean machine. All dependent magnets follow their existing expressions.
Probes use the reference closed-orbit particle as \code{co_guess} and the
shared transfer-matrix and chromatic-step settings. Knobs are restored after
each probe sequence, including when an evaluation fails.

\subsection{Tunes}

The two tune trims are additive offsets to \code{kqf2} and \code{kqd1},
in $\mathrm{m}^{-2}$. Unlike the analytical tune rows in the optics matrix,
\code{build_tune_response} measures their full family response with a forward
probe of $h=10^{-7}\,\mathrm{m}^{-2}$:
\begin{equation}
 R^{Q}_{pb}=\frac{Q_p(\boldsymbol{u}_0+h\boldsymbol{e}_b)
                       -Q_{p,\rm ref}}{h}.
\end{equation}
This is a $2\times2$ matrix, evaluated with \code{chrom=False}.

\subsection{First-order chromaticity}

Define $Q'_p=\dd Q_p/\dd\delta$, returned as \code{dqx} and \code{dqy}.
Four additive sextupole-family offsets control \code{ksf1}, \code{ksf2},
\code{ksd1}, and \code{ksd2}, in $\mathrm{m}^{-3}$.
\code{build_chromaticity_response} returns the $2\times4$ matrix
\begin{equation}
 R^{Q'}_{pb}=\frac{Q'_p(\boldsymbol{u}_0+h\boldsymbol{e}_b)
                       -Q'_p(\boldsymbol{u}_0-h\boldsymbol{e}_b)}{2h},
 \qquad h=10^{-4}\,\mathrm{m}^{-3}\quad\text{by default}.
\end{equation}
Each probe uses \code{chrom=True}; the targets are the shared reference
chromaticities at nominal nonlinear strengths.

\subsection{Second-order chromaticity}

The two knobs \code{dko_arc_f} and \code{dko_arc_d} represent $k_3$ offsets
in $\mathrm{m}^{-4}$. They act on the quadrupoles controlled by the two arc
quadrupole families. For each affected element the installer adds
$\Delta\code{knl[3]}=u_b L_j$ (with its correction switch), preserving the
existing multipole expression. The $2\times2$ response is
\begin{equation}
 R^{Q''}_{pb}=\frac{Q''_p(\boldsymbol{u}_0+h\boldsymbol{e}_b)
                        -Q''_p(\boldsymbol{u}_0-h\boldsymbol{e}_b)}{2h},
 \qquad h=0.01\,\mathrm{m}^{-4}\quad\text{by default}.
\end{equation}
Twiss returns $Q''_p$ as \code{ddqx} and \code{ddqy}. The reported quadratic
coefficient is instead $Q''_p/2$, since
\begin{equation}
 Q_p(\delta)=Q_p(0)+Q'_p\delta+\tfrac12 Q''_p\delta^2+\cdots.
\end{equation}
The optimizer stores derivatives and targets of $Q''$ without dividing by two.
Consequently a tolerance of one on $Q''/2$ corresponds to a target tolerance
of two on \code{ddqx} or \code{ddqy}.

For both chromatic responses, the Twiss momentum step is
\code{delta_chrom=2e-4}. This momentum step is distinct from the family-knob
probe $h$. It is shared by reference and probe calculations; the earlier
$10^{-6}$ momentum step gave unreliable second derivatives in this workflow.
Both chromatic corrections are used in refinement, after the error ramp,
with first-order correction preceding second-order correction.

\section{Implementation map and scope}

The analytical kernels are in \code{response_builders.py}. Assembly and
target ordering are in \code{corr_orbit.py}, \code{corr_optics.py}, and
\code{corr_coupling.py}. The numerical derivatives are in \code{corr_tune.py},
\code{corr_chromaticity.py}, and \code{corr_second_order_chrom.py}.

The unweighted matrices are installed with \code{set_response} by
\code{CorrectionOptimizer}. Evaluations and the correction sequence belong
to \code{FCCeeCorrection}.

The analytical formulas describe a stable periodic reference away from their
resonant denominators. They are not exact responses of a strongly coupled,
misaligned, nonlinear, radiating machine. Finite-difference Twiss tests in
\code{tests/test_fccee_optics_response.py} check normal-optics responses,
including compound knobs and distinct beta/alpha locations. The coupling
directional check provides an additional runtime consistency test. Four-point
quadrature improves the longitudinal integration of the chosen model; it
does not remove that model's approximations. Neither response validation nor
correction convergence establishes dynamic or momentum acceptance.

\appendix
\section{Four-point Gauss--Legendre integration}
\label{app:gauss}

For each quadrupole integral, the code calls
\code{numpy.polynomial.legendre.leggauss(4)}. The returned nodes $\xi_n$ are
the roots of the fourth Legendre polynomial on $[-1,1]$, and $w_n$ are its
quadrature weights:
\begin{center}
\begin{tabular}{rr}
\toprule
Node $\xi_n$ & Weight $w_n$ \\
\midrule
$-0.861136311594053$ & $0.347854845137454$ \\
$-0.339981043584856$ & $0.652145154862546$ \\
$+0.339981043584856$ & $0.652145154862546$ \\
$+0.861136311594053$ & $0.347854845137454$ \\
\bottomrule
\end{tabular}
\end{center}
For a magnet of length $L_j$, map these dimensionless nodes and weights to
physical distances and integration lengths:
\begin{equation}
 \ell_{jn}=\frac{L_j}{2}(1+\xi_n),\qquad
 W_{jn}=\frac{L_j}{2}w_n,\qquad
 \int_0^{L_j} f(\ell)\,\dd\ell
 \simeq\sum_{n=1}^4 W_{jn}f(\ell_{jn}).
 \label{eq:gauss}
\end{equation}
Thus $\sum_nW_{jn}=L_j$, and $W_{jn}$ has units of metres. The rule is exact
for polynomials through degree seven. Our kernels contain transported optics
and trigonometric phase factors, so their integration is still approximate.
No adaptive subdivision or quadrature error estimate is currently performed.

For example, Eq.~\eqref{eq:beta} becomes
\begin{equation}
 B_{\beta_p(i),j}\simeq
 -\frac{\sigma_p\beta_p(s_i)}{2\sin(2\pi Q_p)}
 \sum_{n=1}^4 W_{jn}\beta_p(t_j(\ell_{jn}))
       \cos\Phi_p(i,t_j(\ell_{jn})).
\end{equation}
The same nodes enter the other normal-quadrupole kernels. One evaluates the
\emph{whole response integrand} at each point; averaging beta and phase first
would not give the same result. In particular, beta and phase vary inside a
strong quadrupole, so one point or an entrance/exit average can miss that
variation. Gauss integration retains four samples while avoiding extra Twiss
calculations: local optics are obtained by the transport below.

\section{Optics propagation inside a quadrupole}
\label{app:transport}

\subsection{Constant-gradient transfer matrix}

Let $P=1+\delta_j$ at the quadrupole entrance and
$K_p=\sigma_p k_{1j}/P$. Within the model, $k_{1j}$ and $P$ are constant.
Define
\begin{align}
 C(\ell)&=
 \begin{cases}
 \cos(\sqrt{K_p}\ell), & K_p>0,\\
 1, & K_p=0,\\
 \cosh(\sqrt{-K_p}\ell), & K_p<0,
 \end{cases}\\
 S(\ell)&=
 \begin{cases}
 \sin(\sqrt{K_p}\ell)/\sqrt{K_p}, & K_p>0,\\
 \ell, & K_p=0,\\
 \sinh(\sqrt{-K_p}\ell)/\sqrt{-K_p}, & K_p<0.
 \end{cases}
\end{align}
In canonical transverse coordinates $(q,p_q)$, the matrix is
\begin{equation}
 M_p(\ell)=
 \begin{pmatrix} a&b\\c&d\end{pmatrix}
 =\begin{pmatrix} C&S/P\\-P K_p S&C\end{pmatrix}.
\end{equation}
Its determinant is one. In the zero-gradient limit,
$a=d=1$, $b=\ell/P$, $c=0$, giving a drift. In the nominal on-momentum
reference $P=1$; the helper nevertheless retains the entrance momentum
factor in its local transport. This transport alone does not make the response
kernels general off-momentum formulas.

\subsection{Beta, alpha, phase, and dispersion}

With entrance values $\beta_0,\alpha_0$ and
$\gamma_0=(1+\alpha_0^2)/\beta_0$, covariance transport gives
\begin{align}
 \beta(\ell)&=a^2\beta_0-2ab\alpha_0+b^2\gamma_0,\\
 \alpha(\ell)&=-ac\beta_0+(ad+bc)\alpha_0-bd\gamma_0,\\
 \gamma(\ell)&=c^2\beta_0-2cd\alpha_0+d^2\gamma_0.
\end{align}
The helper only needs $\beta(\ell)$; local alpha is not an input to the
response kernels. The $\alpha(s_i)$ in Eq.~\eqref{eq:alpha} is the value
at the observation point, read directly from the reference table.

The phase increment follows from
\begin{equation}
 \sin\Delta\psi=\frac{b}{\sqrt{\beta_0\beta(\ell)}},\qquad
 \cos\Delta\psi=\frac{a\beta_0-b\alpha_0}
                         {\sqrt{\beta_0\beta(\ell)}}.
\end{equation}
Accordingly, the code uses
\begin{equation}
 \mu(\ell)=\mu_0+
 \frac{\atanTwo(b,a\beta_0-b\alpha_0)}{2\pi}.
\end{equation}
The two-argument arctangent retains the quadrant. The implementation uses
its principal branch without further unwrapping, assuming the modest phase
advance of an individual lattice quadrupole; it is not a general transport
routine for arbitrarily large internal phase advances.

For a straight quadrupole in this clean, uncoupled model, the dispersion is
transported homogeneously:
\begin{equation}
 \begin{pmatrix}D(\ell)\\D_p(\ell)\end{pmatrix}
 =M_p(\ell)\begin{pmatrix}D_0\\D_{p,0}\end{pmatrix},\qquad
 D(\ell)=aD_0+bD_{p,0}.
\end{equation}
$D_p$ is the derivative of canonical momentum with respect to $\delta$
(\code{dpx} or \code{dpy}), not an independently assumed slope coordinate.
There is no bending source inside this transport model. Only $D(\ell)$ is
needed at the integration samples.

The code vectorizes these calculations over all quadrupoles and four samples,
returning arrays of shape $(N_{\rm quad},4)$ for beta, phase, integration
weight, and dispersion. It uses the main \code{k1} field for this internal
transport. The separate knob-to-strength mapping in Eq.~\eqref{eq:chain}
accounts for changes in total integrated normal strength.

\end{document}
